% !TEX root = main.tex
% Shared identities for the three separate step-learning settings in Appendix D.
\providecommand{\LRstep}{L_{\mathrm{step}}}
\providecommand{\LRout}{L_{\mathrm{out}}}
\providecommand{\LRsig}{\sigma}
\providecommand{\LRnorm}[1]{\left\lVert #1\right\rVert}
\providecommand{\LRinner}[2]{\left\langle #1,#2\right\rangle}
\theoremstyle{plain}
\newtheorem{lrlemma}[lemma]{Lemma}
\newtheorem{lrproposition}[proposition]{Proposition}
\newtheorem{lrtheorem}[theorem]{Theorem}
\newtheorem{lrremark}[remark]{Remark}
\crefname{lrlemma}{Lemma}{Lemmas}
\crefname{lrproposition}{Proposition}{Propositions}
\crefname{lrtheorem}{Theorem}{Theorems}
\crefname{lrremark}{Remark}{Remarks}

\subsection{Shared setup and gradient identities}
\label{app:weak-routing-initial}

The step objective and innovation are defined in \cref{def:losses,sec:credit}.
For the binary model let $\chi=p(1-p)$; uniform $s_0$ makes $s_{t-1}$ independent
of the prompt bits. Independence gives
\begin{equation}\label{eq:lr-innovation-cov}
 \mathbb E[(b_t-p)(b_j-p)]=\chi\mathbf1\{j=t\}.
\end{equation}
In the teacher-forced layout,
\begin{equation}\label{eq:lr-constant-lag}
 i_t-j_t=n,\qquad t=2,\ldots,n.
\end{equation}
For arbitrary logits $z_t=Vh_\phi(x_t)$, expanding log-sum-exp minus the target
logit gives the shared identity
\begin{equation}\label{eq:lr-kl-decomp}
 L_{\mathrm{step}}=H(p)+\mathbb E\operatorname{KL}
 (q_{s_{t-1}}\Vert\operatorname{softmax}(z_t))-\Phi_{\mathrm{rule}},\qquad
 \Phi_{\mathrm{rule}}=\mathbb E\langle g_t,z_t\rangle.
\end{equation}
Thus $\nabla_V\Phi_{\mathrm{rule}}=\mathbb E[g_th_t^\top]$; at the exact
state-only predictor the loss readout gradient is its negative.
The supplied-routing proof uses its general-$K$ assumptions; the full-prefix
proof specializes to $K=2,p=1/4$ and full causal normalization.
